\section{Deployment Risks and Mitigations}
\label{sec:risks}

The findings above and the interface documentation point to four deployment risks: 1) text in the state can persuade the model, whether it arrives as an instruction or as evidence; 2) whoever names the options, and possibly whoever orders them, controls part of the decision; 3) a silent model update may change the probability scale on which thresholds were tuned; and 4) schema validity or pooled calibration can be mistaken for correctness on a particular workload. Few defenses have been evaluated against these attacks on typed decision models. \rev{One study of guardrails built on seven open-weight typed models found that a few lines of irrelevant log text, or a misleading name on the permissive option, can open the gate, and that every defense it tested could be defeated. A decision reversed by an attack was no less confident than the original, so a review band based on confidence alone does not catch it~\cite{azizi2026one}.} Most of the steps below are therefore design principles with support from adjacent settings.

\noindent\textbf{Keep fetched content out of the gate (risk 1).} Restricting the gate to the original authorization, validated arguments and known policy reduces its direct exposure to untrusted tool output and retrieved content. This applies the principle of separating control flow from untrusted data to the decision component, as in CaMeL~\cite{debenedetti2025defeating}. The authors of JevAdvBench likewise recommend treating the state as untrusted input~\cite{hu2026jevadvbench}. The cost is utility, because a decision that legitimately depends on tool output needs a separate, narrower question.

\noindent\textbf{Allow, review or block (risks 1 and 4).} We recommend three outcomes for a probabilistic gate, allow, review and block, with human approval for consequential actions. TypeSafe's documentation likewise divides confidence into three ranges with different actions~\cite{typesafe2026docs}. On R-Judge at a 1\% miss budget, Liu reports that separate allow and block thresholds raised Jev's automation coverage from 7.63\% to 47.03\%, mainly through more blocks~\cite{liu2026evaluating}. Liu also finds that thresholds that pass validation need not hold on test data.

\noindent\textbf{Fit thresholds per question (risk 4).} Exploratory analyses link successful attacks to small decision margins. We recommend fitting each threshold on labeled data from the deployment and adding a review band around it. Section~\ref{sec:pilot} gives a rule of thumb for the number of labels.

\noindent\textbf{Fix option names and order yourself (risk 2).} In binary questions, neutral option identifiers largely removed the renaming effect~\cite{sun2026type}. With three or more options, even neutral letters changed about half of Laya's decisions. Sun et~al.\ did not run this test on Jev. The deployer, rather than any party that registers options, should fix the option order. Where the model is open, the readout can be averaged over option permutations~\cite{zheng2024large}. LLM judges can favor outputs labeled with well-known model names~\cite{gu2024survey}, so candidate names should be anonymized when the options are models.

\noindent\textbf{Pin the model version (risk 3).} Thresholds are tuned to one model's probability scale, so a silent update can act as a silent policy change. TypeSafe's \texttt{jev-latest} alias resolves to the most recent stable release, and its documentation advises pinning an explicit version~\cite{typesafe2026docs}.
